\section{Position of the results and mathematical scope}\label{sec:literature}
The conditional-square decomposition and the static quantization identity are standard, and the circle quantizer is classical~\cite{GrafLuschgy,RosenblattRoychowdhury}. The new assertion is not that conditional expectation loses variance. It is the identification of the additional pooled-row obstruction, the exact condition under which every rowwise loss can be attained by one recursively reused labeling, and its consequence for a nonreset relative-report experiment with fully counted deadline memory. The spherical law and the exact balanced circle formula quantify a distinction invisible to the one-time acquired distribution: identical spherical marginals coexist with a horizon-dependent causal cost. The two positive realizations identify the same mechanism without a shared classical eigenbasis or an assumed contracting filter.

The nearest comparison in causal coding is the zero-delay theory for Markov sources. Linder and Y\"uksel~\cite{LinderYuksel} establish structural and existence results for optimal quantizers, including finite-horizon deterministic convex-cell quantizers under their hypotheses. Ghomi, Linder and Y\"uksel~\cite{GhomiLinderYuksel} obtain stationary optimal coding and near-optimal finite-memory codes for linear vector sources. In those models the encoder observes the current source. Here the device sees a relative instrument report and a finite retained label, not the current source point. Observing that point at every call would allow a fresh static quantizer and remove the product obstruction. Our internally charged deadline also differs from a publicly scheduled coding architecture. The conditional expectation tools are shared; neither observation interface is a special notation for the other.

Strong uniform values for partially observed control~\cite{VenelZiliotto} concern one strategy that is near optimal over sufficiently long horizons and a long-run criterion. Finite-memory near-optimality in long-run POMDPs~\cite{ChatterjeeSaonaZiliotto} permits the memory needed for a prescribed accuracy to be chosen. Yu and Bertsekas~\cite{YuBertsekas} establish near optimality of finite-state controllers under the stated average-cost hypotheses. Theorems~\ref{thm:energy} and~\ref{thm:clock-circle} instead fix the complete finite-state resource and the terminal test, quantify all finite horizons, and include the internal deadline. They neither imply a general strong uniform value nor contradict existence of increasingly large nearly optimal controllers. The tracking horizon here is not a claim of decreasing static estimation risk with increasing sample count.

History-space existence and finite-window approximation~\cite{YukselHistory}, filter-contraction average-cost approximation~\cite{DemirciKaraYuksel}, and strategic-measure optimality~\cite{YuStrategic} address substantially broader control or measurability questions. Our raw transport condition is narrower and explicitly verifiable; within it the full-history update is an isometry rather than a strict contraction. General state-dependent report likelihoods are not covered by writing a nonlinear filter as a coordinate map. Multichain policy existence and computation~\cite{Leizarowitz} likewise concern a distinct optimization object from the exact integer allocation solved here. The allocation corollary supplies a genuine arithmetic complexity bound for this class, not a polynomial-time global synthesis theorem for POMDPs.

For the retained physical-time theory, ergodic approximation functionals~\cite{Shaw} and Poisson representations~\cite{GlynnInfanger} are classical ingredients. Its contribution remains the marked duration--exit converse with a maximal-prefix normalization, the indispensable physical-return remainder, the finite signed dual, and the resolution-localized transfer of actual scored occupation. These are three separate layers: a quantitative inequality for a fixed finite marked boundary, a compact-family qualitative/robust theorem under its return hypotheses, and a spatial matching theorem requiring actual-mass witnesses plus implementable covers. Proposition~\ref{prop:access} adds a prescribed-start test transfer; it does not make those hypotheses disappear. The nonreset energy theorem is independent of that renewal proof.

The present compatibility condition is necessary and sufficient for saturating the local energy coefficient in the declared finite test interface. We have not proved that it characterizes every finite-memory causal experiment, that the spherical lower product is attained on every higher-dimensional orbit, or that arbitrary acquired measures admit compatible recursive quantizers. The matching spherical rates do not require that stronger claim. Arbitrary unknown kernels, nondominated common-program compactness, unrestricted sensing controls, unbounded stopping, infinite-horizon optimality, and a fully matched program/workspace/precision/time/minimal-simulator region require additional mathematics. These are specific scope boundaries, not replacements of the general acquired-geometry problem. The results proved above retain that problem's causal and resource content while supplying an exact nonreset criterion and a sharp resource law on a nontrivial raw experiment class.
